\chapter{视频识别、跟踪与检索}
\label{chap:vision-video-recognition}

“刚才把钥匙放在哪里”是Ego4D第一视角经历查询所关注的问题类型，需要把问题对应到过去的画面\scite{vision-ego4d2022}

路口的一辆汽车在十帧中出现，究竟应计作一辆还是十辆？人在杯子旁伸手，究竟是拿起还是放下？一张图可以提供外观与空间关系，回答这些问题却还需要把不同时间的观测联系起来。联系的对象也不同：车辆计数需要身份，拿放判断需要顺序，找到某次拿杯的片段则需要起止时间。

\ref{chap:vision-image-recognition}章已经建立单帧识别、定位和匹配的基础。本章把连续观测组织成可用的时间证据，分别讨论对象判断、身份与范围维护、活动发现及查询检索。各项任务可以共享特征与轨迹，但必须保留自己的输出与评价条件。以下汽车和拿杯序列、分数及代价均为教学构造。

\section{视频任务与可用信息}
\label{sec:vision-video-information}

视频用按时间排列的帧$\bm{X}_1,\ldots,\bm{X}_n$表示，帧的拍摄时刻为$\tau_1,\ldots,\tau_n$。等帧率时，相邻时刻差为$1/f$秒，其中$f$是每秒帧数；可变帧率、丢帧或抽样后，应保存实际时间戳。帧编号相差十，只有在时间采样已知时才能换算成持续时间。

\subsection{输出对象与时间范围}
\label{sec:vision-video-outputs}

同一段拿杯视频可以产生多种标注。“有人拿杯子”是范围内的活动标签；每帧的人物框及共同身份组成框轨迹；杯子的逐帧像素区域组成掩码轨迹；从手接触杯子到杯子离开桌面的区间则描述某种约定下的动作范围。训练时只有视频标签，不能把未提供的逐帧边界当作监督；评价时框位置正确，也不能替代身份或活动正确。

\term{目标跟踪}（object tracking）在连续观测中延续选定对象的位置或区域，使后续结果仍指向同一个实例。给定首帧汽车框，即使不知道它的车型，也可以尝试跟踪；自动发现所有车辆则还需要检测及新目标管理。跟踪返回的身份通常是序列内编号，不是现实世界中的车牌或人物姓名。

裁剪片段通常已经把目标活动置于有限范围，未裁剪视频还含等待、背景和多次重复。若训练只见到以动作为中心的短片段，使用时直接对整小时录像分类，输入与输出条件都发生了变化。镜头切换还可能使图像坐标失去连续性；跨镜头重新认出同一辆车需要额外身份依据。视频时间上的位移描述也不能直接当作带米和秒单位的三维运动。

\subsection{观测范围与信息到达条件}
\label{sec:vision-video-availability}

离线与在线描述信息何时可用，局部片段与长时历史描述处理多少信息。这两个维度可以组合。离线算法可以只看五帧，在线算法也可以保存数小时历史；后者在时刻$\tau$只能使用已经到达的帧。若允许等待$\delta$秒后输出，应把新增后帧和等待时间一起记录，具体范围见图\ref{fig:vision-observation-window}。

经历查询要求知道历史中哪些细节值得保存。Ego4D中的相关任务需要返回过去片段，第一视角的遮挡、相机运动和可见范围又会影响证据是否存在。部分样本有其他传感信息，使用时应逐样本确认，不能统一假设视频都有完整音频或视线记录。

保存每帧全部表示会随时间增长。MovieChat用长短期记忆压缩密集视觉信息，区分全局问题与针对某一时刻的问题\scite{vision-moviechat2024} 压缩后若只剩“有人开车”，后来就难以从这句话恢复车牌字符。Flash-VStream进一步面向持续到达的视频维护概括信息与细节记忆，响应成本要连同帧率、历史长度、硬件和提问时刻解释\scite{vision-flash-vstream2025} 这些系统的语言连接将在\ref{chap:vision-multimodal-models}章说明；此处需要记住，历史存在过不等于回答时仍能访问它。

\section{对象与属性的时序识别}
\label{sec:vision-video-object-recognition}

一辆车穿过画面时，车身、车牌和车灯的最佳观测可能出现在不同帧。时间识别的收益来自这些互补证据。直接把所有帧平均起来，可能把模糊或错配的观测也加进来；必须先确定取什么、是否属于同一对象，再决定怎样汇总。

\subsection{观测与对象证据的选择}
\label{sec:vision-video-evidence-selection}

若相机以每秒30帧记录，一辆车只在0.3秒内经过，用每秒一帧的等间隔抽样可能完全漏掉它。改变起始相位也会改变是否采到车辆。提高采样率增加覆盖，但同时增加检测和编码开销；可以先用较粗的时间索引找出运动或对象出现的范围，再在这些范围内加密读取。这样的筛选仍要检查漏选，不能只评价已找到的片段。

整帧表示适合利用场景上下文，对象裁剪更集中于车身或服饰。裁剪之前若没有可信的框，识别器可能读到路面或另一辆车；相机运动和出画面还会使固定裁剪失效。保存帧号、框坐标和清晰度，才能在识别出错时查明模型看到了什么。车牌需要的分辨率与车型不同，不能仅以“已采到一帧”判断证据充分。

\subsection{跨帧外观证据的聚合}
\label{sec:vision-video-aggregation}

设已经确认若干裁剪属于同一辆车，第$t$次观测对类别$c$给出得分$s_{t,c}$。这些得分可以是分类器的原始输出，也可以是归一化后的类别概率；具体选择应固定。以下用归一化输出演示算术，但并不假定这些数值已经校准。

\subsubsection{逐帧得分与可靠观测的汇总}
\label{sec:vision-video-weighted-scores}

平均汇总让每次观测同等参与，最大值保留最强的单次响应，可靠性加权则让清晰、完整的观测有更大影响。对非负可靠性$r_t$，只要$\sum_t r_t>0$，可写为
\begin{equation}
  \bar{s}_c=\frac{\sum_t r_t s_{t,c}}{\sum_t r_t},
  \qquad
  \widehat{c}=\argmax_c\bar{s}_c\text{。}
  \label{eq:vision-video-weighted-scores}
\end{equation}
分母使权重和为1，避免同一组相对权重整体扩大后改变得分尺度。没有可靠观测时应返回未知或重新取帧，不能把分母为零的情形悄然当作一组零分。

考虑清晰、局部遮挡和严重模糊的三帧。它们对“轿车、货车”的输出分别为$(0.80,0.20)$、$(0.60,0.40)$和$(0.05,0.95)$，可靠性设为$1,0.5,0.1$。均值为$(0.4833,0.5167)$，最大值为$(0.80,0.95)$，二者都选择货车。加权轿车分数却为
\[
  \frac{1\times0.80+0.5\times0.60+0.1\times0.05}{1.6}
  =0.690625\text{，}
\]
货车分数为$0.309375$。这解释了为何应限制模糊误判的影响，但并不证明加权一定改善识别：如果可靠性规则把唯一有效帧误判为不可靠，结果也会变差。

最大值还可能由一次异常高分主导；均值则可能稀释一帧中短暂可见的车牌或标志。人工设定清晰度规则与通过标注学习权重是不同方案。学习方案要交代训练目标、标注来源和验证划分，不能在测试结果上反复挑选有利的权重。

\subsubsection{沿对象轨迹积累属性证据}
\label{sec:vision-video-track-evidence}

跟踪器在每帧返回身份及位置后，区域识别器就能为每个身份单独积累证据。只在关联与可见性可信时，把当前得分写入历史。滑动窗口保留最近$k$次观测，窗外观测完全移除；指数衰减则逐步降低旧证据权重，适合无需保存全部分数的情况。

令$\bm{s}_t$为当前类别得分向量，$r_t$同时表达观测与关联可靠性，$0<\rho<1$是衰减因子。维护分子$\bm{a}_t$和总权重$b_t$：
\begin{equation}
  \bm{a}_t=\rho\bm{a}_{t-1}+r_t\bm{s}_t,\qquad
  b_t=\rho b_{t-1}+r_t,\qquad
  \bar{\bm{s}}_t=\bm{a}_t/b_t\text{。}
  \label{eq:vision-video-decay}
\end{equation}
初始二者都为零，仅在$b_t>0$时输出。无新观测可令$r_t=0$；这时分子、分母同比衰减，旧判断不凭空改变，但应单独标出证据已经陈旧。采样间隔不等时，可按实际时间差设置$\rho=\exp(-\Delta\tau/\kappa)$，其中$\kappa$是以秒计的衰减时间尺度。

以车身颜色为例，过去的红色得分为0.8，新一次可靠观测为0.6，令旧总权重为1、$\rho=0.5$、新权重为1，更新后红色分数为$(0.5\times0.8+0.6)/1.5=2/3$。更新的是统计状态，分类器参数并未重新训练。若遮挡后错误接入另一辆红车，历史积累可能反而加强身份错误；下一节的关联可信性决定这里能否安全汇总。

\subsection{属性与状态的连续判定}
\label{sec:vision-video-continuous-attributes}

车身颜色在短录像中相对稳定，车灯是否点亮却会改变。对变化属性，应输出“哪辆车在何时处于什么状态”，不能用整个轨迹的均值替代状态序列。一个可核验的规则是仅在连续若干次可靠观测都超过阈值时改变状态，并在长时间遮挡后回到未知。阈值和持续长度需要依据事件标注与误报、漏报代价选择。

例如灯光分数连续为$0.2,0.8,0.3,0.9,0.9$，设阈值为0.7、需要连续两次支持，则第三次出现的低分打断先前支持，到第五次才确认点亮。规则能过滤短跳变，也会增加确认延迟，并可能漏掉真实的短闪烁。车灯任务若关心的是转向灯频率，就应保留变化而不是强行平滑。时间识别的目标由用途决定，并非输出越稳定越好。

\section{对象身份与范围的持续维护}
\label{sec:vision-video-tracking}

逐帧汽车分类可以判断画面里有什么，却还不知道当前的框是否延续上一帧的同一辆车。跟踪状态因而不仅包含位置，还要保存身份、运动估计、外观或掩码历史，以及最近一次可靠观测的时间。状态维护要能从初始目标出发，解释一次匹配怎样改变后续判断。

\subsection{目标与初始状态的建立}
\label{sec:vision-video-initialization}

单目标任务可以由用户给首帧框、点或掩码，明确接下来延续哪个目标。多目标任务则需由检测器发现对象，给新出现的合格候选建立编号，并维护多个状态。相同的跟踪算法，在人工准确初始化与自动检测初始化下接收到的证据不同，比较时应记录初始化方式。

\subsubsection{SiamRPN++的模板与搜索区域匹配}
\label{sec:vision-video-siamrpn}

给定首帧汽车框，可以保留包含目标的模板$\bm{Z}$，并在下一帧围绕上次位置裁取较大的搜索区域$\bm{S}_t$。SiamRPN++用共享的深层骨干提取两者特征，随后让模板与搜索特征逐通道相关，把“这里像不像目标”变成空间分布的匹配证据\scite{vision-siamrpn2019}

对某一通道$k$，相关操作可概括为
\begin{equation}
  r_{u,v,k}=\sum_{i,j}z_{i,j,k}\,s_{u+i,v+j,k}\text{。}
  \label{eq:vision-video-depthwise-correlation}
\end{equation}
这里$\bm{Z}$与$\bm{S}_t$在公式中指进入相关计算的特征，$(u,v)$枚举搜索位置。骨干共享之后，模板与搜索分支还经过各自的卷积投影；每个通道先独立比较，再由后续卷积组合通道信息。分类分支为候选锚框预测目标前景得分，回归分支预测相对锚框的中心及尺度偏移。多层特征各形成预测，再以学习的权重融合，使不同层的信息共同决定候选。

训练样本是同一目标的模板和搜索图像对，以锚框与目标真实框的重叠分配前景、背景和回归监督。深层网络中的填充可能带来位置偏好，原方法通过搜索区域的空间移位采样减轻对目标总在中心的依赖。训练结束后，使用时保留首帧模板，逐帧更新搜索区域和目标框；并不每帧重新训练一个类别分类器。

假设汽车A在首帧中心$x=40$像素，搜索范围覆盖下一帧的$x=20$至80像素。候选A在$x=48$处得分0.80，另一辆相似汽车B在$x=70$处得分0.85。实际选框还结合尺度变化惩罚和位置窗口，不能直接把最高相关响应叫作正确身份。这些约束表达连续运动的偏好，但当A突然加速、完全遮挡或离开搜索范围时，同样可能把结果拉向错误位置。原方法固定首帧模板，在线改写模板属于额外设计；跨镜头身份恢复也超出了局部搜索的输入条件。

\subsection{当前观测与历史状态的对应}
\label{sec:vision-video-correspondence}

对应可以发生在像素、局部描述或对象框之间。像素对应帮助解释画面怎样移动，对象对应决定哪个观测延续哪个身份。某个像素的位移估计很精确，并不自动把被遮挡前后的两辆相似车区分开。

\subsubsection{像素与局部区域的对应：RAFT}
\label{sec:vision-video-raft}

\term{光流}（optical flow）是两幅图像间像素的二维位移场。对第一帧的位置$\bm{p}=(x,y)^\top$，位移$\bm{u}(\bm{p})$给出第二帧的取样位置$\bm{p}+\bm{u}(\bm{p})$。横轴向右、纵轴向下，位移以当前网格的像素为单位。这里描述图像中的运动；相机移动、物体运动及透视共同影响它。

RAFT先为两帧提取空间特征，在低分辨率特征网格上计算所有位置对的相关性。设两帧各有$m$个位置、每个位置$d$维，则相关矩阵的元素为
\begin{equation}
  c_{i,j}=\sum_{k=1}^{d}f_{i,k}g_{j,k}\text{。}
  \label{eq:vision-video-all-pairs}
\end{equation}
它为第一帧的每个位置保留第二帧的全局匹配候选。原方法沿第二帧的两个空间维度池化，形成多尺度相关金字塔，而保留第一帧查询位置的分辨率。当前光流估计指定第二帧中的查询中心，再在各尺度的局部邻域读取相关证据\scite{vision-raft2020}

设第$r$轮光流为$\bm{u}^{(r)}$。共享的门控更新模块读取局部相关、第一帧上下文、当前光流及隐状态，预测增量$\Delta\bm{u}^{(r)}$，再更新
\[
  \bm{u}^{(r+1)}=\bm{u}^{(r)}+\Delta\bm{u}^{(r)}\text{。}
\]
相关矩阵是预先构造的，随着估计移动，读取中心改变；隐状态也保留已处理的信息。这与每次重新运行互不联系的匹配器不同。训练对多轮估计施加带衰减权重的光流误差，使用时执行规定的更新轮数，并把粗网格结果上采样到输入分辨率。

例如汽车边缘实际向右移动5个像素，教学更新可以从0开始，依次增加3、1.5和0.5，得到$0,3,4.5,5$。这串数值只演示增量与读取位置的关系，不是RAFT收敛的保证。全相关存储随位置数按$m^2$增长，输入分辨率必须计入资源预算；遮挡位置没有可见的第二帧对应，网络给出的稠密值仍是估计。

要恢复三维速度，还须有深度、相机和时间标定。VGGT-$\Omega$把前馈重建扩展到静态与动态场景，可作为从二维对应转向动态空间恢复的延伸入口\scite{vision-vggt-omega2026} 联合预测仍须说明坐标、尺度和可见性，不能把图像位移直接换成真实世界速度，也不能据单个新模型的局部实验推断所有动态场景都已解决。

\subsubsection{对象与已有轨迹的对应}
\label{sec:vision-video-association}

对象级关联把当前检测框与历史轨迹配对。每条轨迹先预测当前可能位置，再比较候选检测；同一检测不能同时延续两条轨迹，同一轨迹在当前帧也不应接收两个独立目标。因此，逐个选择最相似框并不足够，还要约束配对之间的冲突。

\paragraph{运动预测与一对一分配}
\label{sec:vision-video-motion-assignment}

可用框中心、宽高及它们的速度组成状态$\bm{x}_t$，检测框提供位置与尺寸观测$\bm{z}_t$。线性运动近似下，预测与观测关系为
\[
  \bm{x}_{t|t-1}=\bm{F}\bm{x}_{t-1|t-1},
  \qquad
  \bm{z}_t=\bm{H}\bm{x}_t+\bm{\varepsilon}_t\text{。}
\]
$\bm{F}$按帧间时间推进位置，$\bm{H}$取出框中可观测的量。状态含噪声时，卡尔曼更新利用预测与检测的不确定性调整修正量；此处的作用是把不确定的历史预测和当前测量合并，而非宣布运动一定匀速。

设预测框为$\widehat{\bm{b}}_i$、检测框为$\bm{d}_j$，可取代价$c_{i,j}=1-\operatorname{IoU}(\widehat{\bm{b}}_i,\bm{d}_j)$，排除重叠太小或运动偏离太大的组合。匈牙利算法在允许配对中求总成本最小的一对一分配；实际还要允许未匹配，通过虚拟节点或明确的未匹配代价表达，不能强迫每条轨迹接受一个坏检测。

表\ref{tab:vision-association-cost}给出教学代价。若按行先处理轨迹A，它会占用检测1，B只能接检测2，总成本为$0.10+0.80=0.90$。共同分配则选择A接检测2、B接检测1，总成本为$0.20+0.11=0.31$。检测3在空间门限之外，留下供新目标规则检查。全局分配改善的是这个既定代价目标，代价若不能区分真实身份，最优配对也可能错误。
\begin{table}[htbp]
  \centering
  \begin{tabular}{lccc}
    \toprule
    历史轨迹 & 检测1 & 检测2 & 检测3\\
    \midrule
    A & 0.10 & 0.20 & 不允许\\
    B & 0.11 & 0.80 & 不允许\\
    \bottomrule
  \end{tabular}
  \caption{两条轨迹与三个检测的教学关联代价，越小越合适。“不允许”由配对前的门限排除。例中两个有效匹配均优于约定的未匹配选择，检测3不参加这次分配。}
  \label{tab:vision-association-cost}
\end{table}

外观距离可以加入代价，但需固定表示的训练数据、归一化、权重及门限。强相机运动会使所有框一起移动，急转弯会破坏简单速度外推，两辆相似车交换位置则会同时考验运动和外观。应查看错误来自预测、检测还是关联，而不是把整条轨迹失误都归于检测器。

\paragraph{ByteTrack的高低分检测两次关联}
\label{sec:vision-video-bytetrack}

遮挡后的目标可能仍有一个位置合理但得分较低的检测。若检测后立即把它丢掉，关联阶段就失去这次观测。ByteTrack保留这类候选，将得分高低与是否具有可信历史联系起来\scite{vision-bytetrack2022}

原流程把检测分成高分组与低分组，先将高分组与已有轨迹关联，再用低分组匹配第一次未匹配的可继续跟踪轨迹。第二次主要依据运动预测后的框重叠，避免不可靠外观主导遮挡恢复。未匹配的低分框直接舍弃。作者实现还设置低分下限，只让尚处于跟踪状态的未匹配轨迹参与第二次关联；未匹配的高分框经过未确认轨迹处理及新轨迹阈值，才可能建立身份。已丢失轨迹的保留和重激活另有状态与时间限制，不能把所有历史身份无限制地送入低分关联。

设高分门限为0.6，汽车A、B的当前检测得分分别为0.9和0.3。第一次只匹配A，B暂时未匹配；若低分B框与其预测框重叠充分，第二次就能延续B。同时另一个得分0.2的背景框没有合适历史，会被舍弃。图\ref{fig:vision-bytetrack-association}展示这种条件关系。低分并不等于背景，也不等于一定存在目标，历史用于筛选它的用途。

\begin{figure}[htbp]
  \centering
  % 构造示例：横向位置仅示意关联，不表示实测轨迹。
\begin{tikzpicture}[x=1cm,y=1cm,font=\figurefont\small,
  arr/.style={-{Stealth[length=2mm]},BookInk,line width=0.8pt}]
  \node[anchor=west] at (0,3.3) {上一帧};
  \node[anchor=west] at (3.5,3.3) {当前预测与检测};
  \node[anchor=west] at (7.3,3.3) {关联结果};
  \foreach \y/\name/\col in {2.3/A/BookBlue,0.9/B/BookTeal}{
    \draw[\col,line width=1pt] (0.2,\y-0.35) rectangle (1.6,\y+0.35);
    \node at (0.9,\y) {\name};
    \draw[arr] (1.8,\y) -- (3.4,\y);
    \draw[\col,dashed,line width=0.9pt] (3.6,\y-0.45) rectangle (5.1,\y+0.45);
    \draw[\col,line width=1.1pt] (3.8,\y-0.35) rectangle (5.2,\y+0.35);
    \node at (4.5,\y) {\name};
    \draw[arr] (5.4,\y) -- (7.1,\y);
  }
  \node[above] at (6.2,2.3) {高分0.9};
  \node[above] at (6.2,0.9) {低分0.3};
  \node[anchor=west] at (7.3,2.3) {第一次延续A};
  \node[anchor=west] at (7.3,0.9) {第二次延续B};
  \draw[BookMuted,line width=0.9pt] (3.8,-0.8) rectangle (5.2,-0.1);
  \node at (4.5,-0.45) {背景};
  \draw[arr] (5.4,-0.45) -- (7.1,-0.45);
  \node[above] at (6.2,-0.45) {低分0.2};
  \node[anchor=west] at (7.3,-0.45) {无匹配，舍弃};
\end{tikzpicture}

  \caption{两次关联的教学序列。第一次用高分框延续A，第二次只用位置相符的低分框延续B。虚线框表示预测位置，实线框表示当前检测；右侧孤立低分框没有可信历史而被舍弃，不据此新建身份。}
  \label{fig:vision-bytetrack-association}
\end{figure}

检测器的网络训练与使用时的两次分配是两项过程。关联改变框状态和身份记录，没有在线重新训练检测网络。门限与最大丢失时间需要用验证集确定；如果目标完全没有检测候选，低分阶段也没有新证据可用。持续预测可以等它重新出现，却不能作为已经观测到目标的证明。

\subsection{位置、身份与区域状态的更新}
\label{sec:vision-video-state-update}

框轨迹只给出包围区域，抠取汽车、替换背景或计算精细轮廓还需要逐帧掩码。视频目标分割因此把已知目标的像素信息也保存在历史中，用当前图像查询历史的外观与区域，再恢复当前目标范围。

\subsubsection{XMem的键值读取与分层记忆巩固}
\label{sec:vision-video-xmem}

给定首帧图像和目标掩码，XMem用图像特征构成用于比较的键，用图像与目标掩码共同编码对象值；当前图像的键作为查询。匹配历史位置后，模型把对应的对象值加权读出，结合短时隐状态解码当前掩码\scite{vision-xmem2022}

为看清读取含义，设当前查询位置为$i$、历史位置为$j$，匹配分数为$s_{i,j}$，对象值向量为$\bm{v}_j$。归一化读取写为
\begin{equation}
  a_{i,j}=\frac{\exp s_{i,j}}{\sum_\ell\exp s_{i,\ell}},
  \qquad
  \bm{r}_i=\sum_j a_{i,j}\bm{v}_j\text{。}
  \label{eq:vision-video-memory-read}
\end{equation}
原方法用查询的通道选择权重$e_{i,c}\in[0,1]$和记忆项的缩放系数$\alpha_j\geq1$调节键距离：
\[
  s_{i,j}=-\alpha_j\sum_c e_{i,c}(k_{j,c}-q_{i,c})^2\text{。}
\]
其中$\bm{k}_j$是历史键，$\bm{q}_i$是当前查询。通道选择决定这次比较关注哪些外观维度，缩放系数改变该记忆项在距离增大时衰减多快，两者均由学习模块产生。实际读取还可只保留最高的若干匹配项。若两个候选读取权重为0.8和0.2，值中某个分量分别为0.9和0.1，则该分量读出0.74，再由解码器与当前特征共同形成像素判断，不能直接把0.74当作最终掩码概率。

三类记忆承担不同时间尺度的职责。感知记忆是快速更新的隐状态，保存近期变化；工作记忆保存近期较完整的键值信息；长期记忆用较少原型保存反复有用的历史。工作记忆达到容量后，方法根据读取使用情况选择原型键，并把候选值按与原型的联系汇聚成长期值。长期容量也需管理，以使用频率等信息淘汰部分内容。

汽车被遮挡几秒后重新出现，最近一帧可能只剩遮挡物，长期记忆仍可能保存此前清晰车身。读取结果取决于当前键能否再次指向正确原型，而不是只取决于记忆保存得久。训练通过标注掩码学习编码、读取与解码；使用阶段写入的是特征、预测掩码和隐状态，网络参数保持固定。错误掩码写入会污染后续读取，原型压缩也可能丢掉车牌等细节，有限容量不提供长期身份正确的保证。

\subsubsection{SAM 2的提示分割与流式记忆更新}
\label{sec:vision-video-sam2}

若首帧掩码难以一次标好，可以从点或框提示开始，再针对错误追加提示。SAM 2把这种交互带入视频。本节依据论文2024年10月的修订版，其中更新后的模型对应公开实现中的SAM 2.1。当前图像编码特征通过记忆注意力读取历史空间信息和对象指针；对象指针是从目标解码表示形成的紧凑对象信息，用来参与跨帧条件，而不是现实身份标签\scite{vision-sam2}

提示编码器把点、框或掩码转换为条件，掩码解码器把条件与读过历史的图像特征结合，预测目标区域、质量信息和对象存在得分。不存在或被遮挡时仍要维护状态，但不应据一个看似完整的框认定对象当前可见。预测掩码与图像特征经记忆编码后写入历史。近期非提示帧只保留有限数量，用户提供条件的帧按另外的选择规则参与读取，避免长视频无界存储。

训练使用图像和视频标注，并模拟初始提示及根据预测误差追加的纠正点击，使网络学会利用交互输入。使用时，用户在汽车上点一下，模型传播掩码；若后续车尾漏分，再在漏分区域追加前景点，就会改变该帧的掩码和之后可读取的记忆。纠正影响多远，还取决于实际重新传播了哪些帧，不能只修改当前帧就宣称全部历史已修复。

在线前向传播只能读已到达帧，离线交互可从不同提示时刻向指定方向重新处理给定视频。二者需要不同的观测范围和等待记录。多目标使用时仍需区分各对象的提示、输出及状态；用户点选一辆车后的持续分割，与自动发现并命名所有车辆是不同任务。

\subsection{遮挡、丢失与重新出现的处理}
\label{sec:vision-video-recovery}

可靠匹配后，状态有当前观测支持；暂时失去匹配后，位置来自外推，身份来自历史。输出应能区分这两种来源。等待时间内若出现候选，重新核验空间范围、外观和历史一致性；超过协议期限再终止轨迹，必要时将再次出现的对象作为新身份，而不强行连接。

固定模板避免把错误目标写成新模板，却难以适应外观大变；更新模板能追随姿态与光照，也可能在遮挡时吸收背景。掩码记忆有同样的两难。实际管理可以对低可信观测停止写入、保留早期可靠模板，或允许用户纠正，但这些措施应按具体实现验证。跨镜头相似外观不足以证明同一实例，历史已经淘汰的细节也不能靠延长推理重新获得。

\section{活动与事件的判断}
\label{sec:vision-video-activities}

身份把对象在时间上连起来，活动判断还要解释对象之间发生了什么。“手、杯子和桌子”在拿起与放下时都可能出现，区别在运动方向、接触关系和状态顺序。本节讨论没有外部查询时，从预设活动类别中判断发生了什么，以及是否还要确定发生范围。

\subsection{活动证据的形成}
\label{sec:vision-video-activity-evidence}

外观可以提示场景和对象，光流可以提供局部运动，手与杯子的相对变化则提示交互。证据怎样分布在时间上，决定采样和表示需要保留什么。长动作需要较广的覆盖，短暂运动需要足够密集的观测，两种要求并不相同。

\subsubsection{TSN的分段稀疏采样与类别共识}
\label{sec:vision-video-tsn}

只取相邻几帧可能看到“伸手”，却没有看到杯子之后是否离开桌面。时间分段网络（temporal segment network, TSN）把给定视频均分成$K$段，每段选一个短片段，用共享网络分别输出类别得分，再通过共识函数形成视频判断\scite{vision-tsn2016}

令第$k$段采样片段为$\bm{V}_k$、网络输出为$\bm{s}_k=f_\theta(\bm{V}_k)$。采用均值共识时，先汇总得分，再做softmax及视频标签损失：
\begin{equation}
  \bm{g}=\frac{1}{K}\sum_{k=1}^{K}\bm{s}_k,\qquad
  p_c=\frac{\exp g_c}{\sum_j\exp g_j},\qquad
  L=-\log p_y\text{。}
  \label{eq:vision-video-tsn}
\end{equation}
这里$y$是给定范围的真实活动类别。损失从视频汇总结果反传到共享网络，而不是把视频标签机械当作每一帧都正在发生的活动。训练时在各段随机取样，使用时固定抽帧数量、位置和空间裁剪协议。

在九秒教学视频中，伸手、拿杯和收手分别主要发生在三个时间段。连续抽取开头三帧可能都只见伸手，跨三段取样却有机会覆盖完整过程。图\ref{fig:vision-temporal-sampling}对照了覆盖范围。若三段对“拿杯、放杯”的得分依次为$(1,0)$、$(3,0)$、$(2,0)$，均值为$(2,0)$，拿杯概率约为0.881；交换这三个得分的顺序，均值完全不变。因此，共识汇总本身不编码段间次序。

\begin{figure}[htbp]
  \centering
  % 九秒教学序列，样本线不表示实际方法的固定帧数。
\begin{tikzpicture}[x=0.82cm,y=1cm,font=\figurefont\small]
  \foreach \a/\b/\name in {0/3/伸手,3/6/拿杯,6/9/收手}{
    \fill[BookPaper] (\a,2.45) rectangle (\b,3.05);
    \node at ({(\a+\b)/2},2.75) {\name};
    \draw[BookRule,line width=0.5pt] (\a,0.2) -- (\a,3.1);
  }
  \draw[BookRule,line width=0.5pt] (9,0.2) -- (9,3.1);
  \draw[-{Stealth[length=2mm]},BookInk,line width=0.8pt]
    (0,0) -- (9.3,0);
  \node[anchor=east] at (9.3,-0.65) {时间（秒）};
  \foreach \x in {0,3,6,9}{
    \draw[BookInk,line width=0.7pt] (\x,0.05) -- (\x,-0.05);
    \node[below] at (\x,-0.08) {\x};
  }
  \node[anchor=east] at (-0.2,1.95) {跨段取样};
  \draw[BookMuted,line width=0.6pt] (0,1.95) -- (9,1.95);
  \foreach \x in {1.5,4.5,7.5}{
    \draw[BookBlue,line width=1.5pt] (\x,1.7) -- (\x,2.2);
  }
  \node[anchor=east] at (-0.2,0.95) {连续片段};
  \draw[BookMuted,line width=0.6pt] (0,0.95) -- (9,0.95);
  \fill[BookTeal!12] (3.6,0.65) rectangle (5.4,1.25);
  \foreach \x in {3.7,3.95,4.2,4.45,4.7,4.95,5.2}{
    \draw[BookTeal,line width=1pt] (\x,0.7) -- (\x,1.2);
  }
\end{tikzpicture}

  \caption{稀疏覆盖与连续片段提供的时间证据不同。图中三段和动作位置为教学设定，竖线表示采样帧。跨段采样覆盖伸手、拿杯和收手，但稀疏点间仍可能漏掉短动作；连续片段保留局部运动，却未必覆盖全程。}
  \label{fig:vision-temporal-sampling}
\end{figure}

TSN还可分别以RGB图像和堆叠的光流作为两路输入，各自形成预测后按规定权重融合。光流要由帧间估计取得，其计算成本不应从系统延迟中删除。RGB分支可能依靠杯子和桌子的静态线索，光流分支则多了运动信息；即使融合后活动类别正确，也没有自动给出精确起止边界。

\subsubsection{I3D的连续片段与时空特征}
\label{sec:vision-video-i3d}

要区分一只手正向杯子靠近还是远离，可把连续帧一起输入，让局部滤波同时跨越空间和时间。I3D将图像网络的二维滤波器扩展为三维滤波器，在保留已有图像参数的同时开始学习相邻时刻的关系\scite{vision-i3d2017}

若二维卷积核为$\bm{W}$，把它沿时间复制$q$份，并将每份设为$\bm{W}/q$。对没有运动、每帧相同的输入，远离时间填充边界的卷积响应满足
\begin{equation}
  \sum_{r=1}^{q}\frac{\bm{W}}{q}*\bm{X}
  =\bm{W}*\bm{X}\text{。}
  \label{eq:vision-video-inflation}
\end{equation}
这个初始化使静态重复片段的响应尺度与原图像卷积对应。随后视频训练让不同时间切片的权重分别变化，网络才可能学习先后和运动；“复制”本身尚未学会方向。池化的时间核和步长也要单独设置，不能把所有空间下采样原样照搬到时间维。

I3D从连续RGB或光流片段提取时空特征，通过汇总与分类头输出类别，以视频标签学习，也可利用Kinetics预训练参数适配具体任务。两路输入的来源和计算预算应分别注明。若输入16帧、原帧率30、每隔2帧取一帧，首末观测间隔是$15\times2/30=1$秒；同样16个输入时刻，用不同帧率或步长会看到不同速度和范围。

连续片段保留局部顺序证据，但时间下采样、汇总和有限感受范围仍可能丢掉关键变化。拿杯与放杯在局部看似相反，实际镜头运动、遮挡和背景偏置可能使分类器依赖其他线索。检验运动判断时，可以比较正放、倒放及保持背景而替换动作的受控片段，结果须按具体模型确认。

\subsection{活动类别与发生范围的确定}
\label{sec:vision-video-activity-output}

已经取得时间特征后，输出还可能是一个范围的类别、某个人的动作标签、若干事件区间，或者每个时刻的步骤。网络都在处理视频，并不意味着这些输出可以用同一组标注训练或用同一个分数评价。

\subsubsection{给定时间范围内的活动判定}
\label{sec:vision-video-clip-label}

单标签分类假设一个片段选一个主要类别，通常用softmax及类别交叉熵。若同一人同时站立、拿杯并说话，就需要多个独立标签，可用各类别的sigmoid输出及二元交叉熵训练。阈值决定最终保留哪些标签，类别不均衡与标注遗漏都影响阈值选择。

整段标签还没有说明是谁在做。AVA对关键时刻的人物框标注多种原子动作。原论文方法从关键帧产生人物候选，将相同空间区域在时空特征上汇集，再预测该人的动作标签和框修正，结合人物框与动作监督学习\scite{vision-ava2018} 这一路线把人物区域引入活动判断，但用同一空间范围跨时刻取特征，并不等于已经恢复了精确人物轨迹。

教学画面中，左边的人拿杯、右边的人看手机，两人都站立。整帧输出“拿杯、看手机、站立”没有完成动作到人物的绑定；每个人的多标签才表达所需关系。反过来，人物框和轨迹完全正确也不能替代手物交互的判断。AVA关键时刻的标注规定了评价时刻，不直接给出整次拿杯动作从开始到结束的连续范围。

\subsubsection{事件区间与连续步骤的确定}
\label{sec:vision-video-temporal-boundaries}

\term{时间动作定位}（temporal action localization）在未裁剪视频中输出预设动作类别及发生区间，例如两次拿杯分别位于第9至17秒和第25至30秒。\term{动作分割}（action segmentation）则为每个时间位置分配步骤类别，再把相邻同类位置组成片段，例如拿杯、倒水和放杯。区间检测允许空白、重叠和重复事件，逐帧单标签分段通常要求在规定步骤词表中为每个位置作出选择；并发活动需另设输出。

\paragraph{ActionFormer的多尺度时间点与边界回归}
\label{sec:vision-video-actionformer}

不同事件持续时间差别很大，固定长度窗口难以同时贴合短拿杯和长等待。ActionFormer把按时间排列的片段特征编码为多尺度序列，以局部时间注意力交换邻近信息，并通过下采样扩大每层单个位置对应的时间范围\scite{vision-actionformer2022}

各尺度位置分别预测类别及到左、右边界的非负距离。设某位置换算到物理时间后为$\tau$，该尺度一步对应$\Delta$秒，网络输出距离$a,b$以尺度步为单位，则候选区间为
\begin{equation}
  [\widehat{\tau}_{\mathrm{s}},\widehat{\tau}_{\mathrm{e}}]
  =[\tau-\Delta a,\ \tau+\Delta b]\text{。}
  \label{eq:vision-video-boundaries}
\end{equation}
若输入特征曾按视频时长重新采样，还需先恢复其原始时间映射。距离以秒给出时可令$\Delta=1$；不能把不同尺度的原始输出不加换算地相减。

例如$\tau=12$秒、左右距离已经换算为3秒和5秒，输出为$[9,17]$秒。若模型实际在每步0.5秒的网格上输出$a=6,b=10$，仍得到同一区间。图\ref{fig:vision-temporal-boundaries}用这个例子说明时间点与左右边界的关系。一个位置预测正类不表示只有该瞬间发生动作，它还要通过两侧距离描述持续范围。

\begin{figure}[htbp]
  \centering
  % 上下两部分是独立的教学时间范围；纵向位置不承载数值。
\begin{tikzpicture}[x=1cm,y=1cm,font=\figurefont\small]
  \node[anchor=west] at (0,3.7) {时间点预测左右边界};
  \draw[-{Stealth[length=2mm]},BookInk,line width=0.8pt]
    (1,2.45) -- (10,2.45) node[below left] {秒};
  \fill[BookTeal!14] (2,2.85) rectangle (8.4,3.2);
  \draw[BookTeal,line width=1pt] (2,2.7) -- (2,3.35);
  \draw[BookTeal,line width=1pt] (8.4,2.7) -- (8.4,3.35);
  \fill[BookInk] (4.4,3.02) circle (2pt);
  \draw[{Stealth[length=1.7mm]}-{Stealth[length=1.7mm]},BookInk,line width=0.8pt]
    (2,3.02) -- (4.4,3.02);
  \draw[{Stealth[length=1.7mm]}-{Stealth[length=1.7mm]},BookInk,line width=0.8pt]
    (4.4,3.02) -- (8.4,3.02);
  \node[above] at (3.2,3.1) {3秒};
  \node[above] at (6.4,3.1) {5秒};
  \foreach \x/\lab in {2/9,4.4/12,8.4/17}{
    \draw[BookInk,line width=0.7pt] (\x,2.4) -- (\x,2.5);
    \node[below] at (\x,2.4) {\lab};
  }
  \node[anchor=west] at (0,1.55) {逐帧标签与片段碎裂};
  \node[anchor=east] at (1.2,0.9) {真实};
  \node[anchor=east] at (1.2,0.2) {预测};
  \foreach \x/\gt/\pr in {1.5/拿/拿,2.8/倒/倒,4.1/倒/拿,5.4/倒/倒,6.7/放/放,8/放/放}{
    \draw[BookRule,fill=BookPaper,line width=0.5pt] (\x,0.62) rectangle ++(1.2,0.55);
    \node at (\x+0.6,0.9) {\gt};
    \draw[BookRule,fill=BookPaper,line width=0.5pt] (\x,-0.08) rectangle ++(1.2,0.55);
    \node at (\x+0.6,0.2) {\pr};
  }
  \draw[BookAmber,line width=1.1pt] (4.1,-0.08) rectangle (5.3,0.47);
\end{tikzpicture}

  \caption{时间点回归区间与逐帧步骤分段的教学对照。上半图中12秒位置预测左右距离，形成9至17秒区间；下半图为另一段六个等长时间格的标签，孤立的“拿”把真实连续的“倒”分成两段。两种输出需要分别检查边界与碎片化。}
  \label{fig:vision-temporal-boundaries}
\end{figure}

训练时按照动作区间、各尺度允许的回归长度及中心区域选择正样本，对类别使用焦点损失，对正位置的区间使用距离交并比回归损失。中心采样使正监督集中在动作中心附近，但它不是推理时需要输入的真实边界。使用时解码各尺度候选，筛选分数，并用Soft-NMS降低高度重叠候选的重复影响。

输出类别来自训练词表，不能自动响应任意文字查询。原方法利用给定序列的前后上下文；若事件正在发生，右端点尚未观测到，在线系统必须区分预测的结束位置与已经确认的结束位置。把完整录像的边界成绩搬到即时告警，会改变可用信息条件。

\paragraph{MS-TCN的逐帧预测与多阶段分段修正}
\label{sec:vision-video-mstcn}

对“拿杯—倒水—放杯”的完整过程，逐帧分类可能在倒水中间偶然跳回拿杯，造成过度分段。MS-TCN以多阶段时间卷积修正概率序列：第一阶段读取视频特征，后续阶段读取前一阶段的类别概率，用逐步扩大的时间感受范围判断孤立跳变是否符合上下文\scite{vision-mstcn2019}

每阶段的膨胀时间卷积跳过一定间隔取样，使多层计算覆盖较长范围。原方法使用非因果卷积，同一位置可以读取未来特征，因而多阶段结构不能直接等同于在线处理。训练在各阶段施加逐帧分类损失，并加入截断的相邻对数概率差：
\begin{equation}
  L_{\mathrm{smooth}}
  =\frac{1}{nC}\sum_{t=2}^{n}\sum_{c=1}^{C}
  \min\!\left\{
  \bigl(\log p_{t,c}-\log p_{t-1,c}\bigr)^2,\ \eta^2
  \right\}\text{。}
  \label{eq:vision-video-smoothing}
\end{equation}
这里$n$为序列长度，$C$为类别数，$\eta>0$为截断上限。原文采用上述$nC$归一化；实现需明确序列首项及无效填充位置如何处理，并避免对零概率取对数。原方法计算这项损失时把前一时刻概率当作固定目标，只对当前$p_{t,c}$一侧求梯度。各阶段的分类与平滑损失加权后求和训练整个网络。

在一次教学跳变中，“倒水”概率从0.8降到0.1，对数差绝对值为$\log 8\approx2.079$，未截断平方约为4.324。若示例设$\eta=2$，这一对只贡献4；从0.8到0.7的变化则仅约为0.0178。截断限制大跳变的惩罚，但分类损失仍需要求真实边界两侧取不同标签，平滑不能独自决定哪些跳变应消除。

最终按最后一阶段概率选取各时刻类别，再合并连续同类标签。图中真实标签为“拿、倒、倒、倒、放、放”，预测为“拿、倒、拿、倒、放、放”，六格中只错一格，却多出一个虚假拿杯段和一段碎裂倒水。仅报告帧准确率$5/6$会掩盖步骤结构错误。相反，过强平滑可能删掉真实的短拿杯步骤，仍须检查边界和重复次数。

\subsection{顺序、重复与状态变化的解释}
\label{sec:vision-video-event-relations}

得到事件区间后，可以把每个结果保存为对象身份、类别、开始时间、结束时间及来源位置。两个事件是否相邻、重叠或重复，应由这些记录计算，而不是只从视频总标签推测。“拿起”与“放下”还要结合杯子相对桌面的状态变化；车辆转弯则要结合轨迹方向及道路参照。

两次拿杯若因区间过宽而合并，次数会少算；一次倒水若被分成三个小段，次数会多算。前一段结束恰好早于后一段开始，只提供时间先后，并不证明前者导致后者。更复杂的跨事件语言解释将在\ref{chap:vision-multimodal-understanding}章中连接到答案证据核验。

\section{视频匹配与片段检索}
\label{sec:vision-video-retrieval}

前面从视频出发发现对象与事件，检索则从查询出发决定什么内容值得返回。同一组帧、轨迹和片段表示都可以进入索引，区别在查询要求及相关性的定义。返回一个相似视频只完成了候选筛选，是否包含同一辆车、所需动作及准确时间范围还要分别核验。

\subsection{查询与匹配要求的确定}
\label{sec:vision-video-query}

视频查视频可能寻找相同事件或相似运动，图像查视频可能寻找相似外观，给定车辆照片则可能要求同一实例。即使三种查询都用向量相似度排序，正例定义也不同。任务协议还要规定返回单位是整段视频、身份轨迹还是时间区间，以及边界可以偏差多少。

文字查询“有人把杯子放到桌上”增加了语言与视频的对应。图文表示对齐见\ref{sec:vision-models-contrastive}节；这里先把文字编码器视为能够输出可比较向量的接口，重点追踪它与哪些帧、哪些范围比较。精确语言指代及查询区间定位将在\ref{chap:vision-multimodal-understanding}章展开。

\subsection{候选视频与片段的筛选}
\label{sec:vision-video-candidates}

每段视频一个向量便于快速搜索，但把位置压缩掉后很难回答某个局部证据来自哪里。局部索引保存更多表示及时间信息，成本较大，却能把命中回溯到实际观测。选哪种粒度，取决于查询是否需要身份、运动或精确时间。

\subsubsection{帧、轨迹与局部片段的视觉匹配}
\label{sec:vision-video-visual-matching}

图像查询可以使用\ref{sec:vision-recognition-retrieval}节的描述向量，与逐帧全图或对象裁剪比较。活动查询可与连续片段的表示比较。索引条目至少保留视频编号、时间范围及必要的区域坐标，表示来源还应注明训练与归一化方式；建立索引并没有重新学习表示。

给定查询向量$\bm{q}$和候选帧向量$\bm{v}_t$，帧级最大相似度$\max_t\langle\bm{q},\bm{v}_t\rangle$强调曾经出现过的强匹配；轨迹均值则积累同一对象的外观证据。前者可能被一次错误匹配抬高，后者依赖身份关联。对运动事件，独立帧得分无法区分同样画面元素的不同顺序，需要使用保存局部时间关系的片段表示。

以汽车图片查录像，整帧可能因为同一片道路和天空而命中，裁剪区域却显示车型不同。沿命中位置回查框与相邻帧，才能检验对象对应。若索引窗口为10秒，排序只说明这个10秒窗口值得检查，不能认定动作恰好持续10秒。长历史可先检索粗摘要，再读取细粒度条目，但在线只能回查实际保存且已经到达的信息。

\subsubsection{CLIP4Clip的帧聚合与视频排序}
\label{sec:vision-video-clip4clip}

CLIP4Clip从已有图文编码器出发，把采样视频帧编码成图像向量，文字查询编码成文本向量，再学习适合视频检索的比较方式\scite{vision-clip4clip2021} 最直接的配置对帧向量求均值，再与文本计算余弦相似度；序列配置先用LSTM或带时间位置的Transformer处理帧序列，再汇总。前者的均值本身不保留顺序，后者可通过序列模块把时间关系写入表示。

设一个批次含$B$对匹配视频和文本，$\bm{v}_i,\bm{q}_j$为归一化后的表示，$z_{i,j}$为用于训练的相似度分数。若使用温度$\tau>0$，可令$z_{i,j}=\langle\bm{v}_i,\bm{q}_j\rangle/\tau$。两个方向分别在同一矩阵的行、列中识别正确配对：
\begin{equation}
  \begin{aligned}
    L_{\mathrm{v2t}}&=-\frac{1}{B}\sum_i
      \log\frac{\exp z_{i,i}}{\sum_j\exp z_{i,j}},\\
    L_{\mathrm{t2v}}&=-\frac{1}{B}\sum_i
      \log\frac{\exp z_{i,i}}{\sum_j\exp z_{j,i}}\text{。}
  \end{aligned}
  \label{eq:vision-video-retrieval-loss}
\end{equation}
训练组合两个方向的损失。批内其他配对充当负候选，但如果两段视频实际都符合描述，就可能产生假负例；标注的一一对应并非语义上的唯一匹配。原工作比较冻结范围、学习率和不同聚合器，使用某种配置时应明确哪些预训练参数被微调、哪些新参数从视频文本配对中学习。

使用时可预先编码候选视频，收到查询后计算相似度并排序。教学查询“有人把杯子放到桌上”对三个视频的分数设为0.82、0.70和0.21，依次返回第一、第二、第三个候选。第二个视频若是拿起杯子，较高相似度可能来自相同对象；若第一段有一分钟而放杯只占两秒，全视频平均又可能稀释这一事件。二者分别指向顺序表达与时间粒度问题。

这一流程排序的是事先给定的视频或片段，没有同时回归动作边界。应用若先切成窗口再检索，窗口长度与重叠就是外部条件，不能把窗口端点解释成模型推断的动作起止。候选筛选的表示与精确定位可以组合，评价仍须分别检查。

\subsection{检索结果的核验}
\label{sec:vision-video-retrieval-check}

视频级检索检查相关性和排序，实例查询还要核验对象是否同一，区间查询要核验起止与重复次数。比如正确视频中有两次拿杯，只返回覆盖两次的长片段，可能满足“找一段含拿杯的视频”，却不满足“找到第二次拿杯的时间”。查询语义改变了成功条件。

精细核验可以沿候选范围重新检测、关联和判断状态，再结合查询要求选择结果。预设活动检测依赖固定类别标注，任意文字定位则还依赖语言对齐与指代监督。前者不能仅因输出了区间，就被称作已解决后者。

\section{结果检验与研究方向}
\label{sec:vision-video-evaluation}

视频划分应尽量以源视频、对象或场景为单位，避免把相邻帧分到训练和测试后形成近重复泄漏。比较还要固定抽帧、分辨率、初始化、未来帧可见范围和人工纠正预算。离线算法能读完整事件，在线算法可能尚未看到结尾，两者的差异应先解释输入，再比较输出。

框与掩码需要检查空间重叠，跟踪还检查身份连续性。身份F1（IDF1）在预测身份和真实身份作全局一对一对应后，按满足身份及位置条件的观测数计算
\begin{equation}
  \operatorname{IDF1}
  =\frac{2\,\operatorname{IDTP}}
  {2\,\operatorname{IDTP}+\operatorname{IDFP}+\operatorname{IDFN}}\text{，}
  \label{eq:vision-video-idf1}
\end{equation}
其中$\operatorname{IDTP}$为身份正确的匹配观测数，其余两项分别是身份误报和漏报。若同一真实目标的十次观测被错误拆成两个不同身份、长度为6和4，即使框都准确，全局一对一对应也只能选择一条预测身份，得到$\operatorname{IDTP}=6$、$\operatorname{IDFP}=4$、$\operatorname{IDFN}=4$，IDF1为0.6。这个简化例子说明逐帧检测正确不代表身份维护成功。实际还应报告身份切换、轨迹中断及不同遮挡长度的表现。

事件区间用\term{时间交并比}（temporal intersection over union, tIoU）衡量范围吻合。对两个有效区间$[a,b]$和$[c,d]$，
\begin{equation}
  \operatorname{tIoU}
  =\frac{\max\{0,\min(b,d)-\max(a,c)\}}
  {(b-a)+(d-c)-\max\{0,\min(b,d)-\max(a,c)\}}\text{。}
  \label{eq:vision-video-tiou}
\end{equation}
$[9,17]$与$[10,18]$相交7秒、并集9秒，得分为$7/9$。区间检测的平均精度还要求类别一致、按得分排序并处理同一真实事件的重复命中，宜报告多个tIoU门限。逐帧步骤则同时报告帧准确率与片段级匹配F1或编辑距离，避免长步骤的正确帧掩盖短步骤被删或反复切碎。检索的Recall@$k$和平均精度仍以约定的相关候选为依据。

交通应用可以把这些结果串起来：检测提供汽车范围，关联维持车辆编号，轨迹与道路方向共同支持转弯判断，检索返回需要复核的片段。错误追踪也应沿同一路径进行。天气导致检测置信降低，可能造成漏配；低帧率使运动预测失准，可能引发身份交换；过宽事件区间又可能把两次转弯算成一次。用一个总分类分数难以识别该修哪一环。

长期身份、细粒度活动、低预算观察和未完成事件响应仍受相同约束。增加记忆容量只有在保留了可辨认证据时才有用，缩短延迟则可能减少可见的运动阶段。应把采样频率、有效历史长度、显存、计算时间和等待时间与任务分项结果一起报告，并检查被淘汰的观测会影响哪些未来查询。

\section{本章小结}
\label{sec:vision-video-summary}

十帧中的同一辆车需要通过对应和状态管理归为一个身份，拿起与放下需要借助时间顺序和对象状态区分，找到某次动作则要保留查询所需的片段与边界。更多帧提供了补充证据，也增加了错配、压缩和等待的可能。

对象识别、轨迹维护、事件发现和检索各自保留输出职责，才能把错误定位到具体环节。对每个结果说明实际看过哪些帧、哪些历史仍可访问、哪些位置由预测补出，是把局部模型组合成可核验视频应用的基础。下一章沿着相同的对应关系，讨论如何恢复和传输连续画面本身。
